% Einheit 3 von 4 – Normalform (bank/quadratische-funktionen/e3.jsonl, zusammenbau v0.1)
%% TODO Zweigzeile Teil 1: was hier gelernt wird (Satz in Schülersprache aus dem Einheitstitel) – Einheit 3
\einheitenkopf[e3]{Einheit 3 von 4 · Normalform}
\zweigzeile{ab Kl. 9, je nach Buch bis Kl. 10 · P10 oft}
\setcounter{aufgabe}{14}

%% TODO Titel: Ich-kann-Satz fehlt in der Bank (Platzhalter: Kettenname) – Nr. 15
%% TODO Anweisung: ein Satz über der ersten Teilaufgabe fehlt in der Bank – Nr. 15
\begin{aufgabe}{Normalform}
%% TODO \rechenplatz{n}: Schrittzahl des Grundfalls fehlt in der Bank (nur bei fester Schreibform, 2.3 b)
\begin{teile}
\teil $f(x) = x^2 + 4x - 1$ – welche Form? \kreuz{Scheitelpunktform} \kreuz{Normalform}
\teil $f(x) = x^2 + 3x + 5$ – Schnittpunkt mit der $y$-Achse? ( \leerfeld | \leerfeld )
\teil $f(x) = x^2 + 4x - 2$ – $f(-3)$? f(-3) = \leerfeld
\end{teile}
\begin{gleichungsraster}[2]
\gl{\text{$f(x) = (x - 2)^2 + 5$ – Normalform?}} &
\gl{\text{$f(x) = (x + 4)^2 + 1$ – Normalform?}} \\
\gl{\text{Stimmt $(x + 5)^2 - 8 = x^2 + 10x + 17$? Weise nach oder widerlege.}} &
\end{gleichungsraster}
\begin{teile}
\teil $f(x) = x^2 - 4x + 1$ – lies den Scheitel am Graphen ab.

\begin{ksys}[xmin=-1,xmax=5,ymin=-4,ymax=2,ablesen]
\parabel{1}{2}{-3}{f}
\end{ksys}
S( \leerfeld | \leerfeld )
\teil $f(x) = x^2 - 6x + 5$ – lies den Scheitel ab, stelle die Scheitelpunktform auf und prüfe mit $x = 0$.

\begin{ksys}[xmin=-1,xmax=6,ymin=-5,ymax=1,ablesen]
\parabel{1}{3}{-4}{f}
\end{ksys}
f(x) = \leerfeld
\teil $f(x) = x^2 + 2x - 3$ – wahr oder falsch? \\ $(2|5)$ liegt auf $f$. \kreuz{wahr} \kreuz{falsch} \\ $f$ ist die um $1$ nach links und um $4$ nach unten verschobene Normalparabel. \kreuz{wahr} \kreuz{falsch} \\ Der Schnittpunkt mit der $y$-Achse ist $(-3|0)$. \kreuz{wahr} \kreuz{falsch}

\begin{ksys}[xmin=-4,xmax=3,ymin=-5,ymax=6,ablesen]
\parabel{1}{-1}{-4}{f}
\end{ksys}
\teil Die Abbildung zeigt die Parabel $f(x) = x^2 - 8x + 13$. Gib den Scheitel an und notiere die Scheitelpunktform. \hfill (P10 2025 OS) \\

\begin{ksys}[xmin=-1,xmax=7,ymin=-4,ymax=2,ablesen]
\parabel{1}{4}{-3}{f}
\end{ksys}
S( \leerfeld | \leerfeld ), f(x) = \leerfeld
\end{teile}
\begin{gleichungsraster}[2]
\gl{\text{Weise nach: $(x + 4)^2 - 3 = x^2 + 8x + 13$. (P10 2017 OS)}} & \\
\end{gleichungsraster}
\end{aufgabe}

%% TODO Titel: Ich-kann-Satz fehlt in der Bank (Platzhalter: Kettenname) – Nr. 16
%% TODO Anweisung: ein Satz über der ersten Teilaufgabe fehlt in der Bank – Nr. 16
\begin{aufgabe}{Normalform – Fehler finden, Begründen}
\begin{gleichungsraster}[2]
\gl{\text{Nina rechnet: \rechnung{(x - 4)^2 + 1 &= x^2 - 16 + 1 \\ &= x^2 - 15} Finde den Fehler und korrigiere.}} &
\end{gleichungsraster}
\begin{teile}
\teil Der Graph von $f(x) = x^2 - 7x + 4$ schneidet die $y$-Achse im Punkt $(0|4)$. Warum?
\end{teile}
\end{aufgabe}
